% titel: Blatt 6, A3: Lewis-Formel und Summenformel
% skript: Mortimer Kap. 8 (chemie: Atombindung, Oktettregel)

Die Oktettregel besagt, dass Atome beim Binden die Elektronenkonfiguration
eines Edelgases (8 Valenzelektronen) anstreben. Eine Lewis-Formel zeigt alle
Valenzelektronen (Bindungsstriche und freie Paare); die Summenformel nennt nur
die Atome je Molekül.

\begin{werkzeugkasten}[Oktettregel und ihre Ausnahmen]
\begin{itemize}
  \item \textbf{Erfüllt:} H bekommt 2 (Duett), Atome der 2.\ Periode bekommen 8.
  \item \textbf{Ausnahmen:}
    \begin{itemize}
      \item Elektronenmangel: $\mathrm{Be}$ ($4$), $\mathrm{B}$ ($6$) --- z.\ B.
        $\mathrm{BF_3}$.
      \item Oktetterweiterung $>8$ (periphere $d$-Orbitale ab Periode~3):
        $\mathrm{PCl_5}$, $\mathrm{SF_4}$, $\mathrm{ICl_3}$.
      \item Radikale mit ungerader Elektronenzahl ($\mathrm{NO}$).
    \end{itemize}
\end{itemize}
Die Summenformel zählt die Atome, die Lewis-Struktur zeigt dazu noch die
Anordnung (Zentralatom, Mehrfachbindungen, freie Paare).
\textit{Skript:} Mortimer Kap. 8 (Oktettregel).
\end{werkzeugkasten}

\begin{achtung}[Darauf achten]
\begin{itemize}
  \item Die \textbf{Summenformel} zählt Atome (z.\ B. $\mathrm{C_2H_5NO_2}$); die
        Lewis-Struktur zeigt zusätzlich Bindungen und freie Paare.
  \item Bei Stickstoff ist die Summenformel der typische Verwechslungspunkt:
        ein N-Atom bildet meist 3 Bindungen und trägt 1 freies Paar.
  \item Organische Formeln (Glycin, Ethanol, Aceton): $\mathrm{C}$ ist
        vierbindig, $\mathrm{O}$ zweibindig, $\mathrm{N}$ dreibindig --- die
        Oktettregel hilft, die Struktur richtig zu bauen.
\end{itemize}
\end{achtung}
